\scanpage{113}
-- and when the composites (158) are known, we check that the morphisms $D_J\longrightarrow D_{J'}$ are known,..

We are therefore more particularly interested in the homomorphisms
\[
(164)\qquad D_{d,\delta}\longrightarrow D_\delta\qquad(d>\delta)
\]
which play the role of the ``target morphisms'' whereas the $D_{d,\delta}\unsure{\longrightarrow|-}D_d$ play the role of ``\unsure{source|source} morphisms'', and $D_{d,\delta}$ the role of ``\deleted{\unsure{f}}multiplicity of arrows'' (where inclusion morphisms) -- between elements of sub-multiplicity \inserted{of $X$} associated with $D$.

We have canonical isomorphisms
\[
(165)\qquad D_{d,d'}\times_{D_{d'}}D_{d',d''}\simeq D_{d,d',d''}\qquad\text{if }d>d'>d''
\]
and we deduce \deleted{\unsure{from it}} from this description the canonical morphisms
\diagram{
\((166)\qquad\)
\begin{tikzcd}[ampersand replacement=\&, column sep=large,row sep=small]
\& \& D_d\\
\& D_{d,d'}\arrow[ur]\arrow[dr] \&\\
D_{d,d',d''}\arrow[ur]\arrow[dr,dash] \& \& D_{d'}\\
\& D_{d',d''}\arrow[ur]\arrow[dr] \&\\
\& \& D_{d''}
\end{tikzcd}
}
but not directly \deleted{$D_{d'}$},
\[
(167)\qquad D_{d,d',d''}\longrightarrow D_{d,d''},
\]
but we note that
\[
(168)\qquad D_{d,d''}\longrightarrow D_d\times D_{d''}
\]
\margin{112}
